\chapter{Vectors: Arrows, Lists, and Points}

Every big idea in this book grows out of one small object: the \emph{vector}. Before
we can talk about matrices, spans, projections, eigenvalues, or the singular value
decomposition that quietly powers half of modern machine learning, we have to be
completely comfortable with vectors. Not just able to manipulate them --- comfortable
with them, the way you are comfortable with the number $3$. You can picture $3$, you
can add it, you can scale it, you can spot it in the wild. By the end of this chapter
you should feel the same easy familiarity with a vector like $\colvec{2 \\ 1}$.

Here is the gentle surprise that organizes the whole chapter. A vector is \emph{three
things at once}, and the three pictures always agree:

\begin{itemize}
  \item an \textbf{arrow} --- a directed step, with a length and a heading;
  \item a \textbf{list} of numbers --- something you can type into a computer;
  \item a \textbf{point} --- a location in space, or a single data record.
\end{itemize}

We will look at the same little objects from all three angles until switching between
them becomes automatic. That flexibility is not a side benefit; it \emph{is} the skill.
Geometry gives you intuition, the list of numbers gives you computation, and the point
of view (pun intended) gives you the bridge to data.

\section{An Arrow Is a Step, Not a Place}
\label{ch02:arrow}

Start with the most physical picture. A vector is an \emph{arrow}: it has a length and a
direction, and that is \emph{all} it has. Draw an arrow that goes two units right and
one unit up. That arrow is a vector. It says ``move two right and one up,'' and it says
nothing else.

\begin{intuitionbox}
A vector is a \emph{displacement} --- a recipe for moving. ``Walk three blocks east and
one block north'' is a vector. It does not care where you start. The same instruction
moves you the same way whether you begin at your house or at the library. So the very
same arrow can be drawn anywhere on the page; sliding it around without rotating or
stretching it does not change the vector. Length and direction are everything; position
is nothing.
\end{intuitionbox}

This is the single most common stumbling block at the very start, so let us name it now.

\begin{trapbox}
\textbf{An arrow is not glued to the origin.} A vector has no fixed home. The two arrows
below --- one starting at the origin, one starting somewhere else --- are the \emph{same
vector}, because they have the same length and the same direction. We often \emph{choose}
to draw a vector starting at the origin (we call that the \emph{standard} or
\emph{position} representative), but that is a convenience, not a law.
\end{trapbox}

\begin{figure}[H]
\centering
\begin{tikzpicture}[scale=1.0,>=Stealth]
  \draw[help lines,gray!30] (-0.5,-0.5) grid (6.5,3.5);
  \draw[->,thick,gray!60] (-0.5,0)--(6.7,0) node[right]{$x$};
  \draw[->,thick,gray!60] (0,-0.5)--(0,3.7) node[above]{$y$};
  % standard representative from origin
  \draw[->,very thick,vecblue] (0,0)--(2,1)
       node[midway,above left]{$\vec{v}$};
  % a translated copy
  \draw[->,very thick,vecblue,dash pattern=on 4pt off 2pt] (4,0.5)--(6,1.5)
       node[midway,below right]{same $\vec{v}$};
  \filldraw[vecblue] (0,0) circle (1.6pt);
  \filldraw[vecblue] (4,0.5) circle (1.6pt);
\end{tikzpicture}
\caption{One vector, two drawings. Both arrows go ``2 right, 1 up,'' so both represent
the same vector $\vec{v}$. The solid arrow is the standard representative anchored at the
origin; the dashed copy is the same step taken elsewhere.}
\label{ch02:fig-translate}
\end{figure}

Why insist on this so early? Because the arrow-as-step picture is exactly what makes
vector \emph{addition} obvious in a moment. If a vector were stuck at the origin, adding
two of them would be a mystery. As a step, it is just ``do this move, then do that one.''

\subsection*{Length and direction, made precise}

The \emph{length} (or \emph{magnitude}, or \emph{norm}) of an arrow is how far it
reaches. For the arrow that goes $2$ right and $1$ up, the Pythagorean theorem gives
length $\sqrt{2^2 + 1^2} = \sqrt{5} \approx 2.236$. We write this $\norm{\vec{v}}$. The
\emph{direction} is the heading the arrow points. We will make both of these fully
precise in the chapter on the dot product; for now, hold the picture: \emph{length plus
direction equals arrow equals vector}.

\section{A List Is How You Compute With an Arrow}
\label{ch02:list}

Geometry is wonderful for intuition, but you cannot type an arrow into a computer. So we
pin a coordinate grid onto the plane and \emph{measure} the arrow. The arrow that goes
$2$ right and $1$ up gets recorded as the ordered list of its measurements:
\[
  \vec{v} = \colvec{2 \\ 1}.
\]
The top number, $2$, is how far the arrow reaches in the $x$-direction; the bottom
number, $1$, is how far in the $y$-direction. These numbers are the \emph{components} (or
\emph{coordinates}, or \emph{entries}) of the vector.

\begin{focusbox}
The arrow is the \emph{idea}; the list is the \emph{measurement} of that idea in a chosen
coordinate system. Change the geometry and the list changes with it; the underlying
displacement is what they both point at. We will write vectors as \emph{columns} almost
everywhere in this book, because columns are what matrices act on (you will see why in
the chapter on matrices as machines).
\end{focusbox}

A small but important convention: the \emph{order} matters. The list $\colvec{2 \\ 1}$ is
the arrow ``2 right, 1 up,'' while $\colvec{1 \\ 2}$ is the arrow ``1 right, 2 up.''
Different lists, different arrows. A vector is an \emph{ordered} tuple of numbers.

\begin{remark}
You will see vectors written as columns $\colvec{2 \\ 1}$ and, to save space in running
text, as rows $(2,1)$ or transposed columns $(2,1)\T$. They mean the same arrow. The
distinction between a true \emph{row vector} and a true \emph{column vector} will
start to matter when we multiply by matrices and take transposes; until then, read a row
in prose as ``the column written sideways.''
\end{remark}

\section{Adding Vectors: Walk, Then Walk Again}
\label{ch02:add}

Now the payoff for thinking of a vector as a step. To add two vectors, \emph{take one
step, then take the other from where you landed}. Geometrically: place the tail of the
second arrow at the head of the first; the sum is the new arrow from the very start to
the very end. This is the famous \emph{tip-to-tail} rule.

\begin{figure}[H]
\centering
\begin{tikzpicture}[scale=1.0,>=Stealth]
  \draw[help lines,gray!30] (-0.5,-0.5) grid (5.5,4.5);
  \draw[->,thick,gray!60] (-0.5,0)--(5.7,0) node[right]{$x$};
  \draw[->,thick,gray!60] (0,-0.5)--(0,4.7) node[above]{$y$};
  % u = (3,1)
  \draw[->,very thick,vecblue] (0,0)--(3,1) node[midway,below right]{$\vec{u}$};
  % v = (1,2) starting at head of u
  \draw[->,very thick,vecred] (3,1)--(4,3) node[midway,right]{$\vec{v}$};
  % sum
  \draw[->,very thick,structurepurple] (0,0)--(4,3)
       node[midway,above left]{$\vec{u}+\vec{v}$};
  % faint parallelogram completion
  \draw[->,very thick,vecred,opacity=0.35] (0,0)--(1,2);
  \draw[->,very thick,vecblue,opacity=0.35] (1,2)--(4,3);
  \filldraw[black] (0,0) circle (1.6pt);
\end{tikzpicture}
\caption{Tip-to-tail addition. Start at the origin, walk along $\vec{u}=\colvec{3\\1}$,
then walk along $\vec{v}=\colvec{1\\2}$ from where you landed. The purple arrow from
start to finish is $\vec{u}+\vec{v}=\colvec{4\\3}$. The faint arrows show that adding in
the other order ($\vec{v}$ first) lands in the same place: the two routes complete a
\emph{parallelogram}.}
\label{ch02:fig-add}
\end{figure}

In coordinates this is wonderfully boring, which is exactly what we want: \emph{add the
matching components}.
\[
  \vec{u}+\vec{v}
  = \colvec{3 \\ 1} + \colvec{1 \\ 2}
  = \colvec{3+1 \\ 1+2}
  = \colvec{4 \\ 3}.
\]
The geometric ``walk then walk'' and the algebraic ``add componentwise'' are the same
operation seen from two sides. The total rightward motion is the sum of the two rightward
motions ($3+1$); the total upward motion is the sum of the two upward motions ($1+2$).
That is the entire content of vector addition.

\begin{intuitionbox}
The faint parallelogram in Figure~\ref{ch02:fig-add} is worth a long look. Going
$\vec{u}$ then $\vec{v}$ and going $\vec{v}$ then $\vec{u}$ both land at the opposite
corner. So $\vec{u}+\vec{v}=\vec{v}+\vec{u}$ is not an algebraic accident you have to
memorize --- it is the statement that a parallelogram has two equal opposite corners.
Geometry remembers the rule for you.
\end{intuitionbox}

\begin{example}
\label{ch02:ex-add3}
Three vectors in $\R^3$. Let
\[
  \vec{a}=\colvec{1\\0\\2},\quad
  \vec{b}=\colvec{-1\\3\\1},\quad
  \vec{c}=\colvec{4\\-2\\0}.
\]
Then
\[
  \vec{a}+\vec{b}+\vec{c}
  = \colvec{1+(-1)+4 \\ 0+3+(-2) \\ 2+1+0}
  = \colvec{4 \\ 1 \\ 3}.
\]
Notice we never had to picture three arrows in 3D at once. Componentwise addition does
not care how many dimensions there are, and it does not care how many vectors you chain
together --- just add each coordinate down its own column. This indifference to dimension
is what will let us, two chapters from now, work in $\R^{784}$ (the space of
$28\times 28$ images) without blinking.
\end{example}

\section{Scaling Vectors: Stretch, Shrink, Flip}
\label{ch02:scale}

The second fundamental operation is multiplying a vector by an ordinary number. In this
context a plain number is called a \emph{scalar}, because its job is to \emph{scale} the
arrow --- to change its length without changing the line it lies along.

Geometrically, $2\vec{v}$ is the arrow $\vec{v}$ stretched to twice its length, pointing
the same way. The vector $\tfrac{1}{2}\vec{v}$ is half as long. And $-\vec{v}$ is the same
length but pointing the \emph{opposite} way --- a negative scalar flips the arrow around.

\begin{figure}[H]
\centering
\begin{tikzpicture}[scale=1.0,>=Stealth]
  \draw[help lines,gray!30] (-3.5,-2.5) grid (4.5,2.5);
  \draw[->,thick,gray!60] (-3.5,0)--(4.7,0) node[right]{$x$};
  \draw[->,thick,gray!60] (0,-2.5)--(0,2.7) node[above]{$y$};
  % v = (2,1)
  \draw[->,very thick,vecblue] (0,0)--(2,1) node[above]{$\vec{v}$};
  % 2v
  \draw[->,very thick,structurepurple,opacity=0.8] (0,0)--(4,2)
       node[midway,below right]{$2\vec{v}$};
  % half v
  \draw[->,very thick,targetgreen] (0,0)--(1,0.5) node[below]{$\tfrac12\vec{v}$};
  % -v
  \draw[->,very thick,vecred] (0,0)--(-2,-1) node[below left]{$-\vec{v}$};
\end{tikzpicture}
\caption{Scalar multiples of $\vec{v}=\colvec{2\\1}$ all lie on one line through the
origin. Positive scalars keep the heading and change the length; the negative scalar
$-1$ flips the arrow to point the other way. The whole line is the set of \emph{all}
scalar multiples of $\vec{v}$.}
\label{ch02:fig-scale}
\end{figure}

In coordinates, again, it could not be simpler: \emph{multiply every component by the
scalar}.
\[
  2\colvec{2\\1}=\colvec{4\\2},\qquad
  \tfrac12\colvec{2\\1}=\colvec{1\\0.5},\qquad
  -\colvec{2\\1}=\colvec{-2\\-1}.
\]
Stretching the arrow to double length doubles how far right \emph{and} how far up it
reaches; that is why each component is multiplied by the same $2$.

\begin{focusbox}
The set of all scalar multiples of a single nonzero vector $\vec{v}$ --- that is,
$\set{t\,\vec{v} : t \in \R}$ --- is a \emph{line through the origin}. Sweep $t$ from
$-\infty$ to $+\infty$ and the tip of $t\vec{v}$ traces out the whole line. This is your
first \emph{subspace} (Chapter on vector spaces), and it is the geometric meaning of the
word ``span'' that the next chapter is built around.
\end{focusbox}

\begin{trapbox}
\textbf{Scaling a vector is not the same as moving its tip by a fixed amount.} The vector
$2\vec{v}$ is \emph{not} ``$\vec{v}$ with each component plus $2$.'' Doubling means
\emph{multiply}, not add. From $\colvec{2\\1}$, scaling by $2$ gives $\colvec{4\\2}$,
while adding $2$ to each entry would give $\colvec{4\\3}$ --- a different arrow pointing a
different direction. Multiplication scales along the line; addition (of a constant vector)
slides off it.
\end{trapbox}

\subsection*{Combining the two: linear combinations in disguise}

Addition and scaling are the only two operations vectors come with, and almost everything
in this book is built by using them together. An expression like
\[
  3\vec{u} - 2\vec{v}
\]
means: stretch $\vec{u}$ by $3$, stretch $\vec{v}$ by $2$ and flip it, then add the
results tip-to-tail. With $\vec{u}=\colvec{1\\2}$ and $\vec{v}=\colvec{3\\1}$,
\[
  3\vec{u}-2\vec{v}
  = 3\colvec{1\\2} - 2\colvec{3\\1}
  = \colvec{3\\6} - \colvec{6\\2}
  = \colvec{-3\\4}.
\]
Such an expression is called a \emph{linear combination} of $\vec{u}$ and $\vec{v}$, and
it is so central that the next chapter is devoted entirely to it. For now just notice that
we already have everything we need to build it: scale, then add.

\section{A Point Is Where an Arrow's Tip Lands}
\label{ch02:point}

Here is the third face of a vector, and the one that opens the door to data. Take the
arrow $\colvec{2\\1}$ and draw it in standard position --- tail at the origin. Its tip
lands at a specific spot in the plane: the location two units right and one unit up. That
\emph{location} is the point $(2,1)$.

\begin{intuitionbox}
\textbf{Arrow and point, two sides of one coin.} Anchor a vector at the origin and its
tip names a point. Conversely, any point names the vector that reaches it from the origin.
So $\colvec{2\\1}$ \emph{is} the displacement ``2 right, 1 up'' \emph{and} the location
$(2,1)$, depending on whether you care about the journey or the destination. Same list of
numbers, two readings.
\end{intuitionbox}

This identification is so useful that we usually stop distinguishing them. When we say
``the point $\vec{x}$ in $\R^2$,'' we mean the location named by the vector $\vec{x}$. The
arrow reading is best when you are \emph{moving} (adding displacements, following a
velocity); the point reading is best when you are \emph{locating} (a data record, a pixel,
a position). Train yourself to flip between them at will.

\begin{figure}[H]
\centering
\begin{tikzpicture}[scale=1.0,>=Stealth]
  \draw[help lines,gray!30] (-0.5,-0.5) grid (3.5,2.5);
  \draw[->,thick,gray!60] (-0.5,0)--(3.7,0) node[right]{$x$};
  \draw[->,thick,gray!60] (0,-0.5)--(0,2.7) node[above]{$y$};
  \draw[->,very thick,vecblue] (0,0)--(2,1);
  \filldraw[vecred] (2,1) circle (2.2pt) node[above right]{$(2,1)$};
  \draw[dashed,gray] (2,0)--(2,1);
  \draw[dashed,gray] (0,1)--(2,1);
  \node[below] at (2,0) {$2$};
  \node[left] at (0,1) {$1$};
\end{tikzpicture}
\caption{The arrow $\colvec{2\\1}$ in standard position and the point $(2,1)$ where its
tip lands. The dashed lines read off the components as coordinates. Arrow and point are
the same vector, seen as a journey or as a destination.}
\label{ch02:fig-point}
\end{figure}

\begin{connectionbox}
This is exactly Descartes' great idea, \emph{analytic geometry}: pin a coordinate grid
onto the plane so that every \emph{geometric} point becomes an \emph{algebraic} list of
numbers, and vice versa. Linear algebra is, in a real sense, Descartes' idea taken
seriously and pushed to every dimension at once. When you later study calculus on curves
and surfaces, or solve differential equations whose ``state'' is a vector, you are leaning
on this same arrow $\leftrightarrow$ point dictionary --- the position, velocity, and
acceleration of a moving particle are all vectors anchored at the same moving point.
\end{connectionbox}

\section{The Zero Vector: Standing Still}
\label{ch02:zero}

There is one special vector we must single out: the \emph{zero vector}, written $\vec{0}$.
It is the vector all of whose components are $0$:
\[
  \vec{0} = \colvec{0\\0} \in \R^2,\qquad
  \vec{0} = \colvec{0\\0\\0} \in \R^3,\qquad
  \text{and so on.}
\]
As an arrow it is the ``arrow'' of length zero: the instruction ``do not move.'' As a
point it is the origin. It is the only vector with no direction at all, and it plays the
role for vectors that the number $0$ plays for numbers.

\begin{focusbox}
The zero vector is the \emph{additive identity}: adding it changes nothing,
\[
  \vec{v} + \vec{0} = \vec{v}\qquad\text{for every }\vec{v}.
\]
``Move by $\vec{v}$, then don't move'' lands you exactly where ``move by $\vec{v}$'' did.
And every vector has an \emph{opposite} that undoes it: $\vec{v} + (-\vec{v}) = \vec{0}$,
because walking out along $\vec{v}$ and back along $-\vec{v}$ returns you to start.
\end{focusbox}

\begin{trapbox}
\textbf{The zero vector $\vec{0}$ is not the number $0$.} One is a list of numbers (an
arrow, a point); the other is a single scalar. They are written almost the same and they
cooperate --- $0\,\vec{v}=\vec{0}$ for every $\vec{v}$ --- but they live in different
places. Writing ``$\vec{v}=0$'' when you mean the zero \emph{vector} is a venial sin we
all commit; just know which $0$ you mean. A subtler trap: a subspace must always contain
$\vec{0}$ (Chapter on subspaces). A line through the origin is a subspace; a parallel line
that misses the origin is not, precisely because it has no zero vector.
\end{trapbox}

\section{From the Plane to \texorpdfstring{$\R^n$}{Rn}}
\label{ch02:rn}

Everything so far has lived in the plane, which we call $\R^2$: the set of all lists of
\emph{two} real numbers. The list picture has one decisive advantage: nothing stops at two.

A vector in three-dimensional space $\R^3$ is a list of three numbers,
$\colvec{x\\y\\z}$, an arrow in the room you are sitting in. We can still draw it
(Figure~\ref{ch02:fig-r3}). Then the pictures run out --- nobody can sketch four
perpendicular axes --- but the \emph{lists} sail right on. A vector in $\R^n$ is simply
an ordered list of $n$ real numbers,
\[
  \vec{x} = \colvec{x_1 \\ x_2 \\ \vdots \\ x_n},
  \qquad x_1,\dots,x_n \in \R.
\]
We call $n$ the \emph{dimension} of the space, and $x_i$ the $i$-th \emph{component}.

\begin{figure}[H]
\centering
\begin{tikzpicture}[scale=1.0,>=Stealth]
  % oblique 3D axes: x toward lower-left, y to the right, z up
  \def\xx{-0.62} \def\xy{-0.42}
  \draw[->,thick,gray!60] (0,0)--(3.0,0) node[right]{$y$};
  \draw[->,thick,gray!60] (0,0)--(0,3.0) node[above]{$z$};
  \draw[->,thick,gray!60] (0,0)--({3.0*\xx},{3.0*\xy}) node[below left]{$x$};
  % target vector components: x=2, y=3, z=2  -> point P
  \coordinate (O) at (0,0);
  \coordinate (Px) at ({2*\xx},{2*\xy});          % x part
  \coordinate (Pxy) at ({2*\xx+3},{2*\xy});        % + y part
  \coordinate (P) at ({2*\xx+3},{2*\xy+2});        % + z part
  % component path (faint)
  \draw[->,thick,vecblue,opacity=0.5] (O)--(Px) node[midway,below]{$x$};
  \draw[->,thick,vecred,opacity=0.5] (Px)--(Pxy) node[midway,below]{$y$};
  \draw[->,thick,targetgreen,opacity=0.5] (Pxy)--(P) node[midway,right]{$z$};
  % the vector itself
  \draw[->,very thick,structurepurple] (O)--(P)
       node[midway,above left]{$\vec{x}$};
  \filldraw[structurepurple] (P) circle (1.4pt);
\end{tikzpicture}
\caption{A vector $\vec{x}=\colvec{2\\3\\2}$ in $\R^3$, assembled from its three
components: walk $2$ along the $x$-axis, then $3$ along $y$, then $2$ up the $z$-axis. The
purple arrow from the origin to the corner is the vector; tip-to-tail addition of the
three axis steps still builds it, exactly as in the plane.}
\label{ch02:fig-r3}
\end{figure}

\begin{formalbox}
\begin{definition}[Real $n$-space]
\label{ch02:def-rn}
For a positive integer $n$, \emph{real $n$-space} $\R^n$ is the set of all ordered
$n$-tuples of real numbers,
\[
  \R^n = \set[\Big]{\, \colvec{x_1\\ \vdots \\ x_n} : x_1,\dots,x_n \in \R \,}.
\]
Its elements are \emph{vectors}. Addition and scalar multiplication are defined
\emph{componentwise}: for $\vec{x},\vec{y}\in\R^n$ and a scalar $c\in\R$,
\[
  (\vec{x}+\vec{y})_i = x_i + y_i,
  \qquad
  (c\,\vec{x})_i = c\,x_i,
  \qquad i = 1,\dots,n.
\]
The zero vector $\vec{0}$ has every component equal to $0$.
\end{definition}
\end{formalbox}

These componentwise definitions are not arbitrary --- they are forced on us if we want
the algebra to keep matching the geometry of arrows in every dimension. And they obey a
short list of rules that we will lean on constantly.

\begin{formalbox}
\begin{proposition}[The vector arithmetic rules]
\label{ch02:prop-axioms}
For all $\vec{u},\vec{v},\vec{w}\in\R^n$ and all scalars $c,d\in\R$:
\begin{align*}
  &\text{(commutativity)} & \vec{u}+\vec{v} &= \vec{v}+\vec{u}, \\
  &\text{(associativity)} & (\vec{u}+\vec{v})+\vec{w} &= \vec{u}+(\vec{v}+\vec{w}), \\
  &\text{(identity)} & \vec{v}+\vec{0} &= \vec{v}, \\
  &\text{(inverse)} & \vec{v}+(-\vec{v}) &= \vec{0}, \\
  &\text{(distributivity over vectors)} & c(\vec{u}+\vec{v}) &= c\vec{u}+c\vec{v}, \\
  &\text{(distributivity over scalars)} & (c+d)\vec{v} &= c\vec{v}+d\vec{v}, \\
  &\text{(compatibility)} & c(d\vec{v}) &= (cd)\vec{v}, \\
  &\text{(unit scalar)} & 1\vec{v} &= \vec{v}.
\end{align*}
\end{proposition}
\end{formalbox}

These eight rules are exactly the axioms of a \emph{vector space}, the structure we will
meet in full later. The point to absorb now is how cheaply they are proved: each one is
just an ordinary number fact applied in every slot at once.

\begin{proof}[Proof of commutativity, as a sample]
Compare the two sides component by component. For each $i$,
\[
  (\vec{u}+\vec{v})_i = u_i + v_i = v_i + u_i = (\vec{v}+\vec{u})_i,
\]
where the middle equality is just commutativity of ordinary real-number addition. Since
the two vectors agree in every component, they are equal. The other seven rules fall the
same way: write out one component, use the matching rule for real numbers, and you are
done.
\end{proof}

\begin{aestheticbox}
Pause on what just happened. The geometry of arrows in the flat plane gave us rules ---
parallelograms, stretching, flipping. We translated those rules into componentwise
arithmetic. And that arithmetic makes \emph{perfect sense in $\R^{100}$ or $\R^{1{,}000{,}000}$},
where no arrow can be drawn. So the algebra has carried the geometry past the edge of what
we can see. This is what gives linear algebra its reach: it lets you reason
\emph{geometrically} about spaces far too big to picture, by trusting that the same rules
still hold. When you later hear that a deep network manipulates vectors in a
$4096$-dimensional space, do not panic --- it is the parallelogram law and stretching,
just with more slots.
\end{aestheticbox}

\section{Unit Vectors and the Standard Basis}
\label{ch02:basis}

Some vectors are special enough to earn names. A \emph{unit vector} is any vector of
length exactly $1$. Unit vectors are the natural way to talk about \emph{pure direction}:
strip a vector of its length and what remains is a unit vector pointing the same way.

To turn any nonzero vector into a unit vector, divide it by its own length --- a move
called \emph{normalizing}:
\[
  \hat{\vec{v}} = \frac{\vec{v}}{\norm{\vec{v}}}.
\]
The little hat $\hat{\phantom{v}}$ traditionally marks a unit vector. For example, with
$\vec{v}=\colvec{3\\4}$ we have $\norm{\vec{v}}=\sqrt{3^2+4^2}=\sqrt{25}=5$, so
\[
  \hat{\vec{v}} = \frac{1}{5}\colvec{3\\4} = \colvec{0.6\\0.8},
\]
and you can check $\norm{\hat{\vec{v}}} = \sqrt{0.6^2 + 0.8^2} = \sqrt{0.36+0.64}
= \sqrt{1} = 1$. Same direction as $\vec{v}$, length trimmed to $1$.

\begin{trapbox}
\textbf{You cannot normalize the zero vector.} Its length is $0$, and dividing by $0$ is
undefined --- which matches the geometry, since $\vec{0}$ has no direction to point a unit
vector along. Normalizing is a move reserved for \emph{nonzero} vectors.
\end{trapbox}

Among all unit vectors, a handful are so useful they get reserved symbols. In $\R^n$, the
\emph{standard basis vectors} $\vec{e}_1,\dots,\vec{e}_n$ are the unit vectors that point
straight along each coordinate axis: $\vec{e}_i$ has a $1$ in slot $i$ and $0$ everywhere
else.

\begin{formalbox}
\begin{definition}[Standard basis of $\R^n$]
\label{ch02:def-ei}
The \emph{standard basis vector} $\vec{e}_i \in \R^n$ is the vector whose $i$-th component
is $1$ and whose other components are all $0$. In $\R^3$,
\[
  \vec{e}_1 = \colvec{1\\0\\0},\qquad
  \vec{e}_2 = \colvec{0\\1\\0},\qquad
  \vec{e}_3 = \colvec{0\\0\\1}.
\]
In physics and engineering the $\R^3$ versions are often written
$\hat{\imath},\hat{\jmath},\hat{k}$.
\end{definition}
\end{formalbox}

\begin{figure}[H]
\centering
\begin{tikzpicture}[scale=1.4,>=Stealth]
  \draw[help lines,gray!30] (-0.4,-0.4) grid (2.4,2.4);
  \draw[->,thick,gray!60] (-0.4,0)--(2.6,0) node[right]{$x$};
  \draw[->,thick,gray!60] (0,-0.4)--(0,2.6) node[above]{$y$};
  \draw[->,very thick,vecblue] (0,0)--(1,0) node[below right]{$\vec{e}_1$};
  \draw[->,very thick,vecred] (0,0)--(0,1) node[above left]{$\vec{e}_2$};
  % build (2,1.5) from basis
  \draw[->,thick,vecblue,opacity=0.4] (0,0)--(2,0);
  \draw[->,thick,vecred,opacity=0.4] (2,0)--(2,1.5);
  \draw[->,very thick,structurepurple] (0,0)--(2,1.5)
       node[midway,above left]{$\vec{x}$};
  \filldraw[structurepurple] (2,1.5) circle (1.4pt);
\end{tikzpicture}
\caption{The standard basis $\vec{e}_1,\vec{e}_2$ of $\R^2$ and the vector
$\vec{x}=\colvec{2\\1.5}=2\vec{e}_1+1.5\,\vec{e}_2$. The faint arrows show $\vec{x}$ being
assembled: go $2$ steps along $\vec{e}_1$, then $1.5$ steps along $\vec{e}_2$.}
\label{ch02:fig-basis}
\end{figure}

Here is why these particular vectors matter so much. \emph{Every} vector in $\R^n$ is a
linear combination of the standard basis vectors, and its coordinates \emph{are} the
coefficients:
\[
  \colvec{x_1\\x_2\\\vdots\\x_n}
  = x_1\vec{e}_1 + x_2\vec{e}_2 + \cdots + x_n\vec{e}_n.
\]
For instance $\colvec{2\\1.5} = 2\vec{e}_1 + 1.5\,\vec{e}_2$ in $\R^2$, and
$\colvec{4\\-1\\7} = 4\vec{e}_1 - \vec{e}_2 + 7\vec{e}_3$ in $\R^3$.

\begin{focusbox}
This is the secret meaning of ``a vector is a list of numbers.'' The list
$\colvec{x_1\\\vdots\\x_n}$ is really a set of \emph{instructions}: take $x_1$ steps along
$\vec{e}_1$, $x_2$ steps along $\vec{e}_2$, and so on, then add them up. The standard basis
is the hidden scaffolding behind every coordinate you have ever written. When we get to
\emph{change of basis}, we will swap this scaffolding for a different one --- and the same
arrow will get a different list. The arrow is real; the list depends on the basis you
measure it against.
\end{focusbox}

\begin{labbox}
Verify the assembly by hand in $\R^3$. Let $\vec{x}=\colvec{4\\-1\\7}$. Compute the linear
combination $4\vec{e}_1 - \vec{e}_2 + 7\vec{e}_3$ slot by slot:
\[
  4\colvec{1\\0\\0} - \colvec{0\\1\\0} + 7\colvec{0\\0\\1}
  = \colvec{4\\0\\0} + \colvec{0\\-1\\0} + \colvec{0\\0\\7}
  = \colvec{4\\-1\\7}.
\]
The $\vec{e}_i$ act like labelled buckets: scalar $i$ drops straight into slot $i$ and
nowhere else. That clean separation is exactly what makes coordinates work.
\end{labbox}

\section{Vectors as Data: A Row of Features}
\label{ch02:data}

So far a vector has been a thing in space. But the very same object --- an ordered list of
numbers --- is how the modern world stores information, and this is where linear algebra
stops being geometry class and starts being the language of data.

A single record about something in the world is a list of measured numbers, and that list
\emph{is} a vector. The trick is to let each coordinate axis stand for a different
\emph{feature} rather than a direction in physical space.

\begin{example}[A house as a vector]
\label{ch02:ex-house}
Describe a house by three numbers: its area in square metres, its number of bedrooms, and
its age in years. A particular house then becomes the vector
\[
  \vec{h} = \colvec{120 \\ 3 \\ 15} \in \R^3
  \quad\longleftrightarrow\quad
  \text{(120 m\textsuperscript{2}, \ 3 bedrooms, \ 15 years old)}.
\]
A whole neighbourhood is a \emph{cloud of points} in $\R^3$, one point per house. Suddenly
geometric questions become data questions: ``which two houses are most alike?'' becomes
``which two points are closest?''; ``what's a typical house?'' becomes ``where is the
center of the cloud?'' --- which, we will see, is just the \emph{average} of the vectors,
$\frac{1}{N}(\vec{h}_1 + \cdots + \vec{h}_N)$, vector addition and scaling doing real
work.
\end{example}

Once data is geometry, the operations of this chapter take on plain-English meanings.
Adding vectors and dividing by $N$ \emph{averages} records into a prototype. Scaling a
feature axis \emph{reweights} or \emph{rescales} a feature (multiply the age coordinate by
a constant and you have changed its units). Subtracting one vector from another,
$\vec{h}_1 - \vec{h}_2$, gives the \emph{difference} record --- how the two houses differ,
feature by feature --- and its length will become our measure of how dissimilar they are.

\begin{applicationbox}
\textbf{Word embeddings and the geometry of meaning.} Modern language models begin by
turning every word into a vector --- often in a space like $\R^{768}$ --- chosen so that
words used in similar ways land near each other. The astonishing part is that the
\emph{arrow} picture survives: directions in this space carry meaning, and the
tip-to-tail addition of this chapter does real semantic work. The famous example is
\[
  \vec{v}_{\text{king}} - \vec{v}_{\text{man}} + \vec{v}_{\text{woman}}
  \;\approx\; \vec{v}_{\text{queen}}.
\]
Read it as displacements: the arrow from ``man'' to ``king'' is roughly the same arrow as
from ``woman'' to ``queen'' --- a ``royalty'' displacement --- so adding it to ``woman''
lands you near ``queen.'' That is the parallelogram of Figure~\ref{ch02:fig-add},
operating in $768$ dimensions on the meanings of words. Every recommendation engine,
search ranker, and chatbot you touch is, at bottom, adding and scaling vectors like the
ones in this chapter.
\end{applicationbox}

\begin{strangbox}
\textbf{Before Strang:} Strang opens his course by asking what you can build from a few
vectors using addition and scalar multiplication --- \emph{linear combinations} --- and
he draws the picture before he writes a single equation. This chapter is the warm-up that
makes that opening feel inevitable: if you are at home with arrows, lists, and points, his
first lecture will read like a story you already know the shape of.

\textbf{After Strang:} when you return to this material having seen the whole arc, reread
the data section above. The cloud of feature vectors in Example~\ref{ch02:ex-house} is
exactly the object that the singular value decomposition and principal component analysis
(the late chapters) take apart. Strang's recurring slogan ``column space'' is, right from
the start, a statement about which points your feature vectors can reach by addition and
scaling. It all begins here.
\end{strangbox}

\section{Subtraction and the Arrow Between Two Points}
\label{ch02:subtract}

We close the geometric tour with one more picture that we will use over and over:
subtraction. By definition,
\[
  \vec{u} - \vec{v} = \vec{u} + (-\vec{v}),
\]
so to subtract, flip the second arrow and add. Componentwise, of course, you just subtract
matching entries:
\[
  \colvec{5\\2} - \colvec{1\\3} = \colvec{4\\-1}.
\]
But the picture is the thing to remember. Think of $\vec{u}$ and $\vec{v}$ as
\emph{points} (tips of arrows from the origin). Then $\vec{u}-\vec{v}$ is the arrow that
points \emph{from $\vec{v}$ to $\vec{u}$}: the displacement you would walk to get from the
second point to the first.

\begin{figure}[H]
\centering
\begin{tikzpicture}[scale=1.0,>=Stealth]
  \draw[help lines,gray!30] (-0.5,-0.5) grid (5.5,3.5);
  \draw[->,thick,gray!60] (-0.5,0)--(5.7,0) node[right]{$x$};
  \draw[->,thick,gray!60] (0,-0.5)--(0,3.7) node[above]{$y$};
  % points u and v
  \draw[->,thick,vecblue] (0,0)--(5,2) node[above right]{$\vec{u}$};
  \draw[->,thick,vecred] (0,0)--(1,3) node[above left]{$\vec{v}$};
  % u - v as arrow from v to u
  \draw[->,very thick,structurepurple] (1,3)--(5,2)
       node[midway,above]{$\vec{u}-\vec{v}$};
  \filldraw[vecblue] (5,2) circle (1.8pt);
  \filldraw[vecred] (1,3) circle (1.8pt);
\end{tikzpicture}
\caption{Subtraction as the arrow between points. With $\vec{u}$ and $\vec{v}$ drawn from
the origin, $\vec{u}-\vec{v}$ is the displacement \emph{from the tip of $\vec{v}$ to the
tip of $\vec{u}$}. Check tip-to-tail: $\vec{v} + (\vec{u}-\vec{v}) = \vec{u}$, so this
arrow really does carry you from $\vec{v}$ to $\vec{u}$.}
\label{ch02:fig-subtract}
\end{figure}

This little fact does an enormous amount of work later. The arrow between two data points
is their \emph{difference}, and its length is their \emph{distance} --- the engine behind
nearest-neighbor search, clustering, and the residuals of least-squares fitting. Whenever
you see $\vec{u}-\vec{v}$ in a formula, picture the arrow from $\vec{v}$ to $\vec{u}$ and
the algebra will narrate itself.

\begin{walkbox}
Take this one on a walk. You are standing at point $\vec{v}=\colvec{1\\3}$ and your friend
is at $\vec{u}=\colvec{5\\2}$. \emph{Without computing}, which way do you point to face
them --- and roughly how far is the walk? Now check: the displacement is
$\vec{u}-\vec{v}=\colvec{4\\-1}$ (four blocks toward larger $x$, one block toward smaller
$y$), a walk of length $\sqrt{4^2+(-1)^2}=\sqrt{17}\approx 4.1$ blocks. Did your arm point
the right way before you did any arithmetic? If so, the arrow picture has taken root.
\end{walkbox}

\section{What to Carry Forward}
\label{ch02:summary}

You now hold the same object three ways, and that flexibility is the whole point.

\begin{focusbox}
\textbf{The chapter in six lines.}
\begin{itemize}
  \item A vector is an \emph{arrow} (length and direction), a \emph{list} of numbers
        (its components in a coordinate system), and a \emph{point} (where the arrow's tip
        lands from the origin) --- one object, three readings.
  \item \textbf{Addition} is tip-to-tail / componentwise: $(\vec{u}+\vec{v})_i = u_i+v_i$.
        It is commutative because parallelograms are symmetric.
  \item \textbf{Scaling} stretches, shrinks, or flips along a line through the origin:
        $(c\vec{v})_i = c\,v_i$.
  \item The \textbf{zero vector} $\vec{0}$ is ``don't move,'' the origin, and the additive
        identity --- and it is not the scalar $0$.
  \item $\R^n$ is just lists of $n$ numbers; the geometry of arrows survives into
        dimensions you cannot draw.
  \item The \textbf{standard basis} $\vec{e}_i$ turns coordinates into instructions:
        $\vec{x} = x_1\vec{e}_1 + \cdots + x_n\vec{e}_n$.
\end{itemize}
\end{focusbox}

The next chapter takes the one operation we built everything from --- ``scale, then
add'' --- and asks the question that organizes all of linear algebra: starting from a
handful of vectors, \emph{which other vectors can you reach}? That set of reachable
vectors is the \emph{span}, and chasing it down leads straight to lines, planes,
subspaces, and the four fundamental spaces at the heart of Strang's course. Keep the three
pictures close; we are about to put them to work.

\section{Exercises}
\label{ch02:exercises}

\begin{exercise}[Warm-up arithmetic]
\label{ch02:exr-arith}
Let $\vec{u}=\colvec{2\\-1}$ and $\vec{v}=\colvec{-3\\4}$. Compute, by hand and
componentwise:
(a) $\vec{u}+\vec{v}$; \quad
(b) $\vec{u}-\vec{v}$; \quad
(c) $3\vec{u}$; \quad
(d) $2\vec{u}-\tfrac12\vec{v}$; \quad
(e) $\vec{u}+\vec{v}+\vec{0}$.
\end{exercise}

\begin{exercise}[Draw the parallelogram]
\label{ch02:exr-parallelogram}
With $\vec{u}=\colvec{3\\1}$ and $\vec{v}=\colvec{1\\2}$, draw a coordinate grid and
carefully sketch $\vec{u}$, $\vec{v}$, $\vec{u}+\vec{v}$, and $\vec{u}-\vec{v}$, all from
the origin (draw $\vec{u}-\vec{v}$ also as the arrow between the appropriate tips). Confirm
visually that $\vec{u}$, $\vec{v}$, and $\vec{u}+\vec{v}$ form three corners of a
parallelogram, and identify the fourth corner.
\end{exercise}

\begin{exercise}[Normalizing]
\label{ch02:exr-normalize}
Find the length of each vector and then its unit vector $\hat{\vec{v}}$:
(a) $\colvec{6\\8}$; \quad
(b) $\colvec{1\\1}$; \quad
(c) $\colvec{2\\-1\\2}$ in $\R^3$.
For (b), why does $\hat{\vec{v}}=\colvec{1/\sqrt2 \\ 1/\sqrt2}$ make sense geometrically?
\end{exercise}

\begin{exercise}[Coordinates are coefficients]
\label{ch02:exr-basis}
Write each vector explicitly as a linear combination of standard basis vectors:
(a) $\colvec{5\\-2}$ in $\R^2$; \quad
(b) $\colvec{0\\3\\-7}$ in $\R^3$; \quad
(c) $\colvec{1\\1\\1\\1}$ in $\R^4$.
Then, going the other way, write out the single vector equal to
$2\vec{e}_1 - 5\vec{e}_3$ in $\R^3$.
\end{exercise}

\begin{exercise}[Scaling versus shifting]
\label{ch02:exr-trap}
A classmate claims that ``doubling a vector means adding $2$ to each entry.'' Using
$\vec{v}=\colvec{2\\1}$, compute both the true $2\vec{v}$ and the classmate's
``$\vec{v}+\colvec{2\\2}$,'' draw both results, and explain in one or two sentences the
geometric difference between scaling and shifting.
\end{exercise}

\begin{exercise}[Midpoint and average]
\label{ch02:exr-midpoint}
Two points are $\vec{p}=\colvec{1\\2}$ and $\vec{q}=\colvec{5\\4}$.
(a) Compute the midpoint $\tfrac12(\vec{p}+\vec{q})$ and verify it sits halfway along the
segment from $\vec{p}$ to $\vec{q}$.
(b) Show in general that $\tfrac12(\vec{p}+\vec{q})$ is the midpoint by writing it as
$\vec{p} + \tfrac12(\vec{q}-\vec{p})$ and reading that as ``start at $\vec{p}$, walk
halfway toward $\vec{q}$.''
\end{exercise}

\begin{exercise}[A house cloud]
\label{ch02:exr-data}
Three houses are recorded as $\vec{h}_1=\colvec{100\\2\\30}$,
$\vec{h}_2=\colvec{140\\3\\10}$, $\vec{h}_3=\colvec{120\\4\\20}$ (area, bedrooms, age).
(a) Compute the ``average house'' $\bar{\vec{h}}=\tfrac13(\vec{h}_1+\vec{h}_2+\vec{h}_3)$.
(b) Compute the difference vectors $\vec{h}_2-\vec{h}_1$ and $\vec{h}_3-\vec{h}_1$ and
state, in words, how houses 2 and 3 each differ from house 1.
\end{exercise}

\begin{exercise}[The arrow between points]
\label{ch02:exr-displacement}
For $\vec{u}=\colvec{7\\1}$ and $\vec{v}=\colvec{2\\3}$, find the displacement from
$\vec{v}$ to $\vec{u}$ and the displacement from $\vec{u}$ to $\vec{v}$. How are the two
related, and why does that match the picture of flipping an arrow?
\end{exercise}

\begin{exercise}[Prove a rule]
\label{ch02:exr-proof}
Prove the distributive law $c(\vec{u}+\vec{v}) = c\vec{u} + c\vec{v}$ for all
$\vec{u},\vec{v}\in\R^n$ and all scalars $c\in\R$, by comparing the two sides one
component at a time (imitate the sample proof of commutativity in the text). Which single
fact about ordinary real numbers does your proof rely on?
\end{exercise}

\begin{exercise}[When is a multiple a unit vector?]
\label{ch02:exr-challenge}
Let $\vec{v}=\colvec{3\\4}$. Find \emph{every} scalar $c$ for which $c\vec{v}$ is a unit
vector. (Hint: $\norm{c\vec{v}} = \abs{c}\,\norm{\vec{v}}$; set this equal to $1$.) How
many such $c$ are there, and where do the two resulting vectors point?
\end{exercise}

\begin{exercise}[Collinear or not --- a geometric eye]
\label{ch02:exr-collinear}
Three points are $\vec{a}=\colvec{1\\1}$, $\vec{b}=\colvec{3\\2}$, $\vec{c}=\colvec{7\\4}$.
By comparing the displacement $\vec{b}-\vec{a}$ with $\vec{c}-\vec{a}$, decide whether the
three points lie on a single straight line. (They are collinear exactly when one
displacement is a scalar multiple of the other --- connect this to the line-through-the-origin
picture of scalar multiples.)
\end{exercise}
